\documentclass[11pt]{article}
\usepackage[margin=1in]{geometry}
\usepackage[T1]{fontenc}
\usepackage{graphicx}
\usepackage{hyperref}
\usepackage{enumitem}
\setlength{\parskip}{0.5em}
\setlength{\parindent}{0pt}

\title{transformer attention models for language}
\author{arXiv GraphRAG --- phase-6 deep research}
\date{\today}

\begin{document}
\maketitle

\begin{abstract}
An excellent, highly contemporary review that masterfully connects historical bottlenecks to the most recent architectural shifts and efficiency breakthroughs.
\end{abstract}

\section{The Evolution and Optimization of Attention Mechanisms in Transformer-Based Language Modeling}

The ability to model linguistic structures effectively hinges on a system's capacity to capture complex dependencies across varying scales of text. For decades, the field of natural language processing was defined by the struggle to manage sequential data, eventually leading to the development of foundational sequence modeling and early attention mechanisms designed to mitigate the limitations of recurrent architectures. As language models evolved, the attention mechanism emerged as the central engine for processing information, enabling models to weigh the relevance of different tokens regardless of their distance in a sequence. This shift has fundamentally redefined how machines interpret context, making the refinement of attention the primary driver of progress in modern artificial intelligence.

This review traces the trajectory of attention-based modeling, beginning with the paradigm shift brought about by the rise of pre-trained transformers and large-scale models. We examine how these models were refined through specific architectural enhancements and positional encoding techniques to better capture the nuances of syntax and order. However, as model scales increased, the quadratic computational cost necessitated a focus on efficiency and complexity reduction in attention. This evolution has been driven by the dual necessity of modeling long-range dependencies and context while maintaining computational feasibility. Furthermore, we explore the efforts to move beyond black-box implementations by analyzing and interpreting transformer mechanisms. Finally, the review concludes by surveying the current frontier, where research is pivoting toward hybrid, sparse, and forgetful attention mechanisms to overcome the inherent constraints of traditional dense attention.

\subsection{Foundations of Sequence Modeling and Early Attention}

Early sequence modeling was defined by recurrent architectures designed to map input sequences to target sequences. The RNN Encoder-Decoder model (\texttt{arXiv:1406.1078}) established a foundational framework for statistical machine translation by utilizing two recurrent neural networks (RNNs): one to encode a source sequence into a fixed-length vector and another to decode that vector into a target sequence. This approach was further specialized in \texttt{arXiv:1409.3215}, which employed a multilayered Long Short-Term Memory (LSTM) architecture to perform end-to-end sequence-to-sequence learning for English-to-French translation.

Despite their utility, these early recurrent models were constrained by a "fixed-length vector bottleneck," where the entire meaning of a source sentence had to be compressed into a single vector. To overcome this, \texttt{arXiv:1409.0473} introduced a soft-search attention mechanism that extended the basic encoder-decoder architecture. This allowed the model to automatically search for relevant parts of a source sentence during the translation process, rather than relying on a static representation. The nuances of this mechanism were further refined in \texttt{arXiv:1508.04025}, which compared a global approach—attending to all source words—against a local approach that focuses on a specific subset of source words to improve neural machine translation performance.

The utility of attention mechanisms eventually expanded beyond text. In \texttt{arXiv:1502.03044}, the concept of attention was applied to image captioning, where the model learned to "fix its gaze" on salient visual objects while generating descriptive words. This demonstrated that the ability to selectively focus on specific inputs was a versatile principle applicable to both linguistic and visual domains.

A fundamental shift in the field occurred with the introduction of the Transformer (\texttt{arXiv:1706.03762}). While previous models, such as those in \texttt{arXiv:1406.1078} and \texttt{arXiv:1409.3215}, were built upon the sequential processing of recurrence, the Transformer's primary novelty lay in its total elimination of recurrence and convolutions. By relying solely on multi-head self-attention, the architecture achieved significantly higher parallelizability and reduced training times compared to its predecessors. This departure from recurrent foundations proved highly effective, as the Transformer established new state-of-the-art BLEU scores, specifically 28.4 for English-to-German and 41.8 for English-to-French translation, marking a transition from recurrent-based sequence modeling to purely attention-based architectures.

\subsection{The Rise of Pre-trained Transformers and Large-Scale Models}

The evolution of transformer-based language modeling has transitioned from establishing bidirectional context to massive-scale parameter expansion. BERT (Bidirectional Encoder Representations from Transformers) established a foundational paradigm by jointly conditioning on both left and right context in all layers through Masked Language Model (MLM) and Next Sentence Prediction (NSP) objectives [\texttt{arXiv:1810.04805}]. This approach achieved significant state-of-the-art results, including an 80.5\% GLUE score and 93.2 SQuAD v1.1 Test F1 [\texttt{arXiv:1810.04805}]. However, BERT's methodology introduced a structural mismatch between pre-training and fine-tuning due to the use of [MASK] tokens, and its larger variant, BERTLARGE, exhibited instability on smaller datasets [\texttt{arXiv:1810.04805}].

Subsequent research addressed these limitations through varying architectural and optimization strategies. RoBERTa sought to improve upon BERT not through structural changes, but through robust optimization of hyperparameters and increased training data size [\texttt{arXiv:1907.11692}]. In contrast, XLNet attempted to resolve the pre-train-finetune discrepancy by utilizing a generalized autoregressive pretraining method [\texttt{arXiv:1906.08237}]. By maximizing expected likelihood over all permutations of factorization order and integrating Transformer-XL concepts, XLNet enabled bidirectional context learning without the [MASK] token issue [\texttt{arXiv:1906.08237}]. While BERT and XLNet focused on encoder-based understanding, BART introduced a denoising sequence-to-sequence architecture, combining a bidirectional encoder with a left-to-right decoder to support generation, translation, and comprehension tasks [\texttt{arXiv:1910.13461}]. Efficiency also became a primary driver of innovation, as seen in ALBERT, which utilized two parameter-reduction techniques and a specialized inter-sentence coherence loss to decrease memory consumption and increase training speed [\texttt{arXiv:1909.11942}].

The trajectory of the field eventually shifted toward extreme scaling. GPT-3 moved away from the bidirectional encoder focus of BERT and XLNet, instead scaling the model to 175 billion parameters to achieve task-agnostic, few-shot performance across domains such as reasoning and text generation [\texttt{arXiv:2005.14165}]. While GPT-3 demonstrated the power of scale, it has faced critiques regarding its performance [\texttt{arXiv:2303.08774}]. This scaling trend culminated in GPT-4, a transformer-based model designed to predict the next token, which expanded the application domain from text-only to multimodal inputs (image and text) [\texttt{arXiv:2303.08774}]. Together, these developments illustrate a move from optimizing specific bidirectional attention mechanisms to leveraging massive parameter counts for generalized, multimodal intelligence.

\subsection{Architectural Enhancements and Positional Encoding}

A central challenge in transformer architectures is the effective integration of sequence order and context, as standard self-attention lacks inherent awareness of position. Early advancements focused on moving beyond absolute position representations to capture the distance between elements. For instance, \texttt{arXiv:1803.02155} introduced a relation-aware self-attention mechanism that utilizes relative position representations to efficiently consider the distances between sequence elements, specifically targeting machine translation tasks. This approach was further refined by \texttt{arXiv:1901.02860}, which introduced Transformer-XL. Transformer-XL addresses context fragmentation by implementing a segment-level recurrence mechanism and a novel positional encoding scheme, allowing the model to learn dependencies beyond a fixed length. This architectural shift significantly improves evaluation efficiency, performing up to 1,800+ times faster than vanilla Transformers.

While these methods focus on relative distance, \texttt{arXiv:2104.09864} proposes RoFormer, which utilizes Rotary Position Embedding (RoPE). RoFormer offers a unique synthesis by encoding absolute positions via a rotation matrix, which simultaneously incorporates explicit relative position dependency directly into the self-attention mechanism. This provides a more mathematically integrated way to leverage positional information compared to the additive or recurrent methods seen in \texttt{arXiv:1803.02155} and \texttt{arXiv:1901.02860}.

Beyond spatial positioning, recent research has expanded the definition of "context" to include temporal and structural dimensions. \texttt{arXiv:2202.02093} introduces temporal attention, a time-aware mechanism designed to capture time-specific contextualized word representations. Unlike the distance-focused models, this architecture is optimized for temporal generalization, making it particularly effective for semantic change detection across different eras in languages such as English, German, and Latin.

Parallel to these refinements in attention mechanisms, other researchers have sought to improve computational efficiency or structural modeling. \texttt{arXiv:1901.10430} proposes dynamic convolutions as a lightweight alternative to self-attention. By predicting separate convolution kernels based solely on the current time-step, this model achieves linear computational complexity relative to input length, offering simplicity for large-scale machine translation and summarization. Finally, to address the inherent limitations of transformers in modeling context-free language tasks, \texttt{arXiv:2405.04515} introduces a differentiable, stack-based attention mechanism. This augmentation moves away from standard attention patterns to incorporate stack-based logic, providing a structural enhancement specifically for natural language tasks that require formal grammatical modeling. Together, these developments illustrate a transition from simple positional encoding to complex, multi-dimensional attention mechanisms that account for relative distance, temporal shifts, and structural hierarchy.

\subsection{Efficiency and Complexity Reduction in Attention}

The primary bottleneck in standard Transformer architectures is the quadratic growth of time and memory requirements relative to sequence length. To address this, several strategies have emerged, ranging from sparse factorizations to hardware-aware exact attention.

Initial efforts to mitigate this scaling issue include the Sparse Transformer (\texttt{arXiv:1904.10509}), which utilizes sparse factorizations of the attention matrix to reduce complexity to \$O(n \textbackslash{}sqrt\{n\})\$ [\texttt{arXiv:1904.10509}]. Subsequent research has pushed toward even greater efficiency. The Reformer (\texttt{arXiv:2001.04451}) achieves \$O(L \textbackslash{}log L)\$ complexity by implementing locality-sensitive hashing (LSH) attention and incorporating reversible residual layers to optimize memory and speed [\texttt{arXiv:2001.04451}].

A significant portion of the literature focuses on achieving true linear complexity. Linformer (\texttt{arXiv:2006.04768}) achieves \$O(n)\$ complexity in both time and space by approximating self-attention with a low-rank matrix [\texttt{arXiv:2006.04768}]. In contrast, BigBird (\texttt{arXiv:2007.14062}) attains linear complexity through a sparse attention mechanism that includes \$O(1)\$ global tokens designed to attend to the entire sequence [\texttt{arXiv:2007.14062}]. While Linformer relies on low-rank approximation, Performers (\texttt{arXiv:2009.14794}) utilize the FAVOR+ approach—using positive Orthogonal Random features—to approximate softmax attention-kernels, aiming to maintain accuracy while achieving linear space and time complexity [\texttt{arXiv:2009.14794}].

A distinct paradigm shift is observed in FlashAttention (\texttt{arXiv:2205.14135}), which moves away from approximating the attention mechanism itself. Instead, it focuses on IO-awareness to improve speed and memory usage for exact attention [\texttt{arXiv:2205.14135}]. By employing tiling, FlashAttention reduces the frequency of memory reads and writes between GPU HBM and on-chip SRAM [\texttt{arXiv:2205.14135}]. This distinguishes it from the approximation-based models like Linformer or Performers, as it maintains the integrity of the exact attention mechanism while optimizing hardware utilization.

The evolution of these models highlights a fundamental trade-off between mathematical approximation and hardware efficiency. While models like Performers (\texttt{arXiv:2009.14794}) and FlashAttention (\texttt{arXiv:2205.14135}) have provided significant advancements in efficiency, they have also been subject to critical analysis regarding their implementation or theoretical guarantees [\texttt{arXiv:2506.11305}]. Ultimately, the field has progressed from reducing complexity via \$O(n \textbackslash{}sqrt\{n\})\$ factorizations [\texttt{arXiv:1904.10509}] and \$O(L \textbackslash{}log L)\$ hashing [\texttt{arXiv:2001.04451}] toward \$O(n)\$ approximations [\texttt{arXiv:2006.04768}, \texttt{arXiv:2009.14794}] and IO-aware exact implementations [\texttt{arXiv:2205.14135}].

\subsection{Modeling Long-Range Dependencies and Context}

A central challenge in transformer architectures is overcoming the quadratic scaling of standard self-attention to facilitate the processing of extended sequences. Several approaches attempt to mitigate this complexity by restructuring how context is captured, ranging from convolutional approximations to memory compression and hybrid attention mechanisms.

One strategy involves replacing self-attention with more efficient operations. The dynamic convolutions model (\texttt{arXiv:1901.10430}) optimizes for efficiency and simplicity by predicting separate convolution kernels based solely on the current time-step to determine context importance. This approach achieves linear computational complexity relative to input length, making it suitable for large-scale machine translation and abstractive summarization. While this offers a lightweight alternative to self-attention, it differs fundamentally from the hybrid attention mechanisms seen in later models. For instance, the Longformer (\texttt{arXiv:2004.05150}) addresses quadratic scaling by combining local windowed attention with task-motivated global attention. Unlike the purely convolutional approach of \texttt{arXiv:1901.10430}, the Longformer serves as a drop-in replacement for standard self-attention, enabling linear scaling for long-document tasks such as WikiHop and generative sequence-to-sequence tasks.

Other models focus on how historical information is retained. The Compressive Transformer (\texttt{arXiv:1911.05507}) optimizes for long-range sequence learning by compressing past memories, a method applied across language modeling (e.g., WikiText-103, PG-19) and reinforcement learning. This memory-centric approach contrasts with the Long-Short Transformer (\texttt{arXiv:2107.02192}), which achieves linear complexity through a dual attention mechanism: long-range attention with dynamic projection paired with short-term attention. The Long-Short Transformer introduces a unique architectural innovation via a dual normalization strategy designed specifically to resolve the scale mismatch between long-range and short-term attention components, a nuance not explicitly addressed by the compression-based strategy of \texttt{arXiv:1911.05507}.

Finally, the modeling of context can be extended beyond linear sequences to structural representations. The Graph-to-Graph Transformer (\texttt{arXiv:2310.17936}) moves away from traditional sequence-based attention by treating attention weights as functionally equivalent to graph edges. By integrating input graph edges directly into attention weight computations and employing iterative graph refinement, this model provides a joint embedding of input, output, and latent graphs. This allows for the modeling of complex linguistic structures through non-autoregressive graph prediction, representing a significant departure from the sequence-length optimization goals shared by the Longformer (\texttt{arXiv:2004.05150}) and the Compressive Transformer (\texttt{arXiv:1911.05507}).

\subsection{Analyzing and Interpreting Transformer Mechanisms}

Research into transformer architectures has transitioned from optimizing structural efficiency to deeply interrogating the functional role of attention heads and their contribution to emergent capabilities. Early investigations into the necessity of multi-headed attention (MHA) focused on architectural redundancy. \texttt{arXiv:1905.10650} demonstrates that the high number of attention heads in models like BERT and various machine translation systems is not strictly necessary for performance, proposing that multiple heads can be pruned at test time to improve speed, memory efficiency, and accuracy. This suggests a level of functional overlap within the MHA mechanism that allows for significant compression without catastrophic loss in linguistic modeling.

While \texttt{arXiv:1905.10650} views heads through the lens of structural redundancy, \texttt{arXiv:1906.04341} shifts the focus toward semantic interpretation, specifically analyzing what BERT’s attention mechanisms learn about language. This work seeks to understand the qualitative aspects of language captured by pre-trained models, moving beyond mere efficiency to probe the underlying representational logic. This interpretive trajectory is furthered by \texttt{arXiv:2305.03210}, which introduces AttentionViz to provide a global view of transformer attention. Unlike previous localized analyses, AttentionViz utilizes a joint embedding of query and key vectors to visualize global patterns across both language and vision transformers, offering a more holistic method for model understanding.

A more specialized functional analysis is provided by \texttt{arXiv:2209.11895}, which identifies "induction heads" as a specific mechanism responsible for in-context learning. By observing how loss decreases at increasing token indices, this work provides a mechanistic explanation for how transformers perform pattern matching, a significant step up from the general semantic observations in \texttt{arXiv:1906.04341}. This focus on specific functional units contrasts with the broader structural pruning discussed in \texttt{arXiv:1905.10650}; where the former seeks to remove heads to save resources, the latter seeks to identify the specific heads that drive complex reasoning.

Finally, the literature explores the limits of these mechanisms regarding data retention. \texttt{arXiv:2411.10115} investigates memorization capacity within attention-only transformers, distinguishing between exact and approximate memorization. This adds a layer of complexity to the understanding of attention: while induction heads (\texttt{arXiv:2209.11895}) facilitate learning through pattern completion, the attention mechanism must also balance the capacity to memorize specific training instances. Together, these papers illustrate a progression from optimizing the quantity of attention heads (\texttt{arXiv:1905.10650}) to visualizing their global interactions (\texttt{arXiv:2305.03210}) and finally decoding the specific mechanistic drivers of learning and memorization (\texttt{arXiv:2209.11895}; \texttt{arXiv:2411.10115}).

\subsection{The Current Frontier: Hybrid, Sparse, and Forgetful Attention}

Recent advancements in attention mechanisms reflect a divergence from standard dot-product attention toward architectures that prioritize computational efficiency, structural locality, or selective information retention. While traditional Transformers struggle with scaling, new models attempt to bridge the gap between the parallel training capabilities of Transformers and the constant-memory inference of Recurrent Neural Networks (RNNs).

A primary direction involves hybridizing attention with other structural priors to improve sequence modeling. The Convolution-Augmented Transformer (CAT) integrates convolutional filters directly into the Key, Query, and Value (K/Q/V) embeddings, aiming to synergize local inductive biases with global attention to enhance recall and length generalization in language modeling \texttt{arXiv:2407.05591}. In contrast, Aaren approaches this hybridity by reformulating attention as a many-to-many RNN through a parallel prefix scan algorithm \texttt{arXiv:2405.13956}. While CAT focuses on embedding-level locality, Aaren focuses on the operational efficiency of the update mechanism, achieving constant memory during inference—a significant advantage for time series forecasting and reinforcement learning where efficient token updates are critical \texttt{arXiv:2405.13956}.

Another frontier explores the necessity of learned attention weights themselves. MixiT challenges the fundamental role of trainable attention by utilizing entirely random attention scores, a design that maintains provably stable signal propagation to disentangle the contributions of attention from the MLP component in language modeling \texttt{arXiv:2506.01115}. This investigation into the necessity of attention is further expanded by the systematic deconstruction of dot-product attention in \texttt{arXiv:2510.11602}, which tests various relaxations of mixing and sequence dependency to identify the core design principles required for effective language modeling.

The challenge of long-context modeling has also led to "forgetful" or selective architectures. The Forgetting Transformer (FoX) introduces a data-dependent forget gate that down-weights unnormalized attention scores, specifically targeting improvements in length extrapolation and short-context downstream tasks \texttt{arXiv:2503.02130}. This mechanism of selective focus is taken to an extreme in the Avey model, which seeks to decouple sequence length from context width entirely \texttt{arXiv:2506.11305}. Rather than modifying the attention score, Avey pairs a ranker with an autoregressive neural processor to select only the most relevant tokens for contextualization, representing a radical departure from both standard attention and traditional recurrence to handle arbitrarily long sequences \texttt{arXiv:2506.11305}. Collectively, these models suggest that the future of language modeling may lie not in denser attention, but in more specialized, selective, or structurally constrained mechanisms.

\subsection{Tensions and Open Problems}

A fundamental tension persists regarding the optimal pre-training objective for capturing linguistic nuance. While early breakthroughs like ELMo (\texttt{arXiv:1802.05365}) established the value of deep contextualized word representations through bidirectional language modeling, subsequent research has questioned whether such approaches sufficiently capture dependency structures. For instance, XLNet (\texttt{arXiv:1906.08237}) argues that traditional autoregressive or bidirectional models fail to account for permutation-based dependencies, proposing a generalized autoregressive framework to overcome the limitations of standard masked language modeling. This debate is further complicated by the shift toward sequence-to-sequence architectures; BART (\texttt{arXiv:1910.13461}) suggests that denoising autoencoding provides a more robust foundation for natural language generation than the pure contextualized representations championed by ELMo (\texttt{arXiv:1802.05365}).

Another critical divide concerns the necessity and efficiency of the attention mechanism itself. A significant body of work has focused on optimizing the quadratic complexity of self-attention, with methods like Linformer (\texttt{arXiv:2006.04768}) advocating for linear complexity to improve upon standard distillation techniques (\texttt{arXiv:1503.02531}). Similarly, sparse computation strategies such as Combiner (\texttt{arXiv:2107.05768}) attempt to maintain transformer efficacy while reducing overhead. However, this "efficiency-first" paradigm is facing increasing scrutiny. Recent critiques, most notably in "Don't Pay Attention" (\texttt{arXiv:2506.11305}), argue that many popular approximations—including the kernel-based methods in Performers (\texttt{arXiv:2009.14794}), the recurrent formulations of fast autoregressive transformers (\texttt{arXiv:2006.16236}), and even hardware-optimized exact attention like FlashAttention (\texttt{arXiv:2205.14135})—may be fundamentally misaligned with the structural requirements of complex sequence modeling. This suggests a growing skepticism toward whether attention-based scaling is the most principled path forward.

Finally, there is an ongoing disagreement regarding the relationship between scale and emergent capabilities. The paradigm of few-shot learning popularized by GPT-3 (\texttt{arXiv:2005.14165}) posits that increasing parameter count and data volume leads to predictable improvements in task generalization. Yet, the technical documentation for GPT-4 (\texttt{arXiv:2303.08774}) provides a more nuanced critique, suggesting that the leap in capability is not merely a linear function of scale but involves complex interactions between instruction tuning and reasoning depth that few-shot scaling laws do not fully capture. This tension is echoed in the debate over scaling laws themselves, as PaLM (\texttt{arXiv:2204.02311}) demonstrates that architectural pathways and training stability can significantly alter the performance trajectories predicted by earlier, more generalized scaling models (\texttt{arXiv:2001.08361}).

\subsubsection{Outlook}

The field is currently moving away from pure attention-scaling toward a synthesis of structural inductive biases and computational efficiency. As we look toward the next generation of models, the frontier lies in reconciling the high-capacity reasoning of transformers with more efficient, perhaps non-attention-based, architectures. Emerging research into convolution-augmented transformers (\texttt{arXiv:2407.05591}) and stack attention (\texttt{arXiv:2405.04515}) suggests that the future may not belong to "pure" attention, but to hybrid models that integrate local structural priors with global relational reasoning, potentially moving toward a more mathematically grounded understanding of how sequence models actually internalize language.

\section*{References}
\begin{enumerate}\small
  \item One Billion Word Benchmark for Measuring Progress in Statistical
  Language Modeling (2013). \texttt{arXiv:1312.3005}
  \item Learning Phrase Representations using RNN Encoder-Decoder for
  Statistical Machine Translation (2014). \texttt{arXiv:1406.1078}
  \item Neural Machine Translation by Jointly Learning to Align and Translate (2014). \texttt{arXiv:1409.0473}
  \item Sequence to Sequence Learning with Neural Networks (2014). \texttt{arXiv:1409.3215}
  \item Adam: A Method for Stochastic Optimization (2014). \texttt{arXiv:1412.6980}
  \item Neural Turing Machines (2014). \texttt{arXiv:1410.5401}
  \item Neural Machine Translation of Rare Words with Subword Units (2015). \texttt{arXiv:1508.07909}
  \item Deep Residual Learning for Image Recognition (2015). \texttt{arXiv:1512.03385}
  \item Aligning Books and Movies: Towards Story-like Visual Explanations by
  Watching Movies and Reading Books (2015). \texttt{arXiv:1506.06724}
  \item Effective Approaches to Attention-based Neural Machine Translation (2015). \texttt{arXiv:1508.04025}
  \item Distilling the Knowledge in a Neural Network (2015). \texttt{arXiv:1503.02531}
  \item Show, Attend and Tell: Neural Image Caption Generation with Visual
  Attention (2015). \texttt{arXiv:1502.03044}
  \item Layer Normalization (2016). \texttt{arXiv:1607.06450}
  \item SQuAD: 100,000+ Questions for Machine Comprehension of Text (2016). \texttt{arXiv:1606.05250}
  \item A Decomposable Attention Model for Natural Language Inference (2016). \texttt{arXiv:1606.01933}
  \item Google's Neural Machine Translation System: Bridging the Gap between
  Human and Machine Translation (2016). \texttt{arXiv:1609.08144}
  \item Gaussian Error Linear Units (GELUs) (2016). \texttt{arXiv:1606.08415}
  \item Pointer Sentinel Mixture Models (2016). \texttt{arXiv:1609.07843}
  \item Using the Output Embedding to Improve Language Models (2016). \texttt{arXiv:1608.05859}
  \item Language Modeling with Gated Convolutional Networks (2016). \texttt{arXiv:1612.08083}
  \item Can Active Memory Replace Attention? (2016). \texttt{arXiv:1610.08613}
  \item Attention Is All You Need (2017). \texttt{arXiv:1706.03762}
  \item Convolutional Sequence to Sequence Learning (2017). \texttt{arXiv:1705.03122}
  \item Outrageously Large Neural Networks: The Sparsely-Gated
  Mixture-of-Experts Layer (2017). \texttt{arXiv:1701.06538}
  \item A Broad-Coverage Challenge Corpus for Sentence Understanding through
  Inference (2017). \texttt{arXiv:1704.05426}
  \item Decoupled Weight Decay Regularization (2017). \texttt{arXiv:1711.05101}
  \item BERT: Pre-training of Deep Bidirectional Transformers for Language
  Understanding (2018). \texttt{arXiv:1810.04805}
  \item Deep contextualized word representations (2018). \texttt{arXiv:1802.05365}
  \item GLUE: A Multi-Task Benchmark and Analysis Platform for Natural Language
  Understanding (2018). \texttt{arXiv:1804.07461}
  \item Character-Level Language Modeling with Deeper Self-Attention (2018). \texttt{arXiv:1808.04444}
  \item Self-Attention with Relative Position Representations (2018). \texttt{arXiv:1803.02155}
  \item Scaling Neural Machine Translation (2018). \texttt{arXiv:1806.00187}
  \item Adaptive Input Representations for Neural Language Modeling (2018). \texttt{arXiv:1809.10853}
  \item Know What You Don't Know: Unanswerable Questions for SQuAD (2018). \texttt{arXiv:1806.03822}
  \item SentencePiece: A simple and language independent subword tokenizer and
  detokenizer for Neural Text Processing (2018). \texttt{arXiv:1808.06226}
  \item The Best of Both Worlds: Combining Recent Advances in Neural Machine
  Translation (2018). \texttt{arXiv:1804.09849}
  \item Universal Transformers (2018). \texttt{arXiv:1807.03819}
  \item Think you have Solved Question Answering? Try ARC, the AI2 Reasoning
  Challenge (2018). \texttt{arXiv:1803.05457}
  \item RoBERTa: A Robustly Optimized BERT Pretraining Approach (2019). \texttt{arXiv:1907.11692}
  \item Transformer-XL: Attentive Language Models Beyond a Fixed-Length Context (2019). \texttt{arXiv:1901.02860}
  \item Generating Long Sequences with Sparse Transformers (2019). \texttt{arXiv:1904.10509}
  \item XLNet: Generalized Autoregressive Pretraining for Language Understanding (2019). \texttt{arXiv:1906.08237}
  \item Exploring the Limits of Transfer Learning with a Unified Text-to-Text
  Transformer (2019). \texttt{arXiv:1910.10683}
  \item ALBERT: A Lite BERT for Self-supervised Learning of Language
  Representations (2019). \texttt{arXiv:1909.11942}
  \item Analyzing Multi-Head Self-Attention: Specialized Heads Do the Heavy
  Lifting, the Rest Can Be Pruned (2019). \texttt{arXiv:1905.09418}
  \item Are Sixteen Heads Really Better than One? (2019). \texttt{arXiv:1905.10650}
  \item HuggingFace's Transformers: State-of-the-art Natural Language Processing (2019). \texttt{arXiv:1910.03771}
  \item What Does BERT Look At? An Analysis of BERT's Attention (2019). \texttt{arXiv:1906.04341}
  \item Adaptive Attention Span in Transformers (2019). \texttt{arXiv:1905.07799}
  \item Compressive Transformers for Long-Range Sequence Modelling (2019). \texttt{arXiv:1911.05507}
  \item Attention is not Explanation (2019). \texttt{arXiv:1902.10186}
  \item DistilBERT, a distilled version of BERT: smaller, faster, cheaper and
  lighter (2019). \texttt{arXiv:1910.01108}
  \item SuperGLUE: A Stickier Benchmark for General-Purpose Language
  Understanding Systems (2019). \texttt{arXiv:1905.00537}
  \item Pay Less Attention with Lightweight and Dynamic Convolutions (2019). \texttt{arXiv:1901.10430}
  \item SpanBERT: Improving Pre-training by Representing and Predicting Spans (2019). \texttt{arXiv:1907.10529}
  \item Megatron-LM: Training Multi-Billion Parameter Language Models Using
  Model Parallelism (2019). \texttt{arXiv:1909.08053}
  \item Are Transformers universal approximators of sequence-to-sequence
  functions? (2019). \texttt{arXiv:1912.10077}
  \item fairseq: A Fast, Extensible Toolkit for Sequence Modeling (2019). \texttt{arXiv:1904.01038}
  \item BERT Rediscovers the Classical NLP Pipeline (2019). \texttt{arXiv:1905.05950}
  \item BART: Denoising Sequence-to-Sequence Pre-training for Natural Language
  Generation, Translation, and Comprehension (2019). \texttt{arXiv:1910.13461}
  \item Theoretical Limitations of Self-Attention in Neural Sequence Models (2019). \texttt{arXiv:1906.06755}
  \item TinyBERT: Distilling BERT for Natural Language Understanding (2019). \texttt{arXiv:1909.10351}
  \item HellaSwag: Can a Machine Really Finish Your Sentence? (2019). \texttt{arXiv:1905.07830}
  \item PyTorch: An Imperative Style, High-Performance Deep Learning Library (2019). \texttt{arXiv:1912.01703}
  \item WinoGrande: An Adversarial Winograd Schema Challenge at Scale (2019). \texttt{arXiv:1907.10641}
  \item Attention is not not Explanation (2019). \texttt{arXiv:1908.04626}
  \item Revealing the Dark Secrets of BERT (2019). \texttt{arXiv:1908.08593}
  \item Transformer++ (2020). \texttt{arXiv:2003.04974}
  \item Learning Hard Retrieval Decoder Attention for Transformers (2020). \texttt{arXiv:2009.14658}
  \item Fixed Encoder Self-Attention Patterns in Transformer-Based Machine
  Translation (2020). \texttt{arXiv:2002.10260}
  \item Guiding Attention for Self-Supervised Learning with Transformers (2020). \texttt{arXiv:2010.02399}
  \item Language Models are Few-Shot Learners (2020). \texttt{arXiv:2005.14165}
  \item Longformer: The Long-Document Transformer (2020). \texttt{arXiv:2004.05150}
  \item Reformer: The Efficient Transformer (2020). \texttt{arXiv:2001.04451}
  \item An Image is Worth 16x16 Words: Transformers for Image Recognition at
  Scale (2020). \texttt{arXiv:2010.11929}
  \item Linformer: Self-Attention with Linear Complexity (2020). \texttt{arXiv:2006.04768}
  \item Transformers are RNNs: Fast Autoregressive Transformers with Linear
  Attention (2020). \texttt{arXiv:2006.16236}
  \item Big Bird: Transformers for Longer Sequences (2020). \texttt{arXiv:2007.14062}
  \item Rethinking Attention with Performers (2020). \texttt{arXiv:2009.14794}
  \item Scaling Laws for Neural Language Models (2020). \texttt{arXiv:2001.08361}
  \item Efficient Transformers: A Survey (2020). \texttt{arXiv:2009.06732}
  \item Long Range Arena: A Benchmark for Efficient Transformers (2020). \texttt{arXiv:2011.04006}
  \item Efficient Content-Based Sparse Attention with Routing Transformers (2020). \texttt{arXiv:2003.05997}
  \item ETC: Encoding Long and Structured Inputs in Transformers (2020). \texttt{arXiv:2004.08483}
  \item ELECTRA: Pre-training Text Encoders as Discriminators Rather Than
  Generators (2020). \texttt{arXiv:2003.10555}
  \item GroupBERT: Enhanced Transformer Architecture with Efficient Grouped
  Structures (2021). \texttt{arXiv:2106.05822}
  \item Long-Short Transformer: Efficient Transformers for Language and Vision (2021). \texttt{arXiv:2107.02192}
  \item TransfoRNN: Capturing the Sequential Information in Self-Attention
  Representations for Language Modeling (2021). \texttt{arXiv:2104.01572}
  \item Evolving Attention with Residual Convolutions (2021). \texttt{arXiv:2102.12895}
  \item Combiner: Full Attention Transformer with Sparse Computation Cost (2021). \texttt{arXiv:2107.05768}
  \item RoFormer: Enhanced Transformer with Rotary Position Embedding (2021). \texttt{arXiv:2104.09864}
  \item Multiformer: A Head-Configurable Transformer-Based Model for Direct
  Speech Translation (2022). \texttt{arXiv:2205.07100}
  \item Relaxed Attention for Transformer Models (2022). \texttt{arXiv:2209.09735}
  \item Temporal Attention for Language Models (2022). \texttt{arXiv:2202.02093}
  \item Pre-Training a Graph Recurrent Network for Language Representation (2022). \texttt{arXiv:2209.03834}
  \item Parallel Attention Forcing for Machine Translation (2022). \texttt{arXiv:2211.03237}
  \item PaLM: Scaling Language Modeling with Pathways (2022). \texttt{arXiv:2204.02311}
  \item In-context Learning and Induction Heads (2022). \texttt{arXiv:2209.11895}
  \item FlashAttention: Fast and Memory-Efficient Exact Attention with
  IO-Awareness (2022). \texttt{arXiv:2205.14135}
  \item What Can Transformers Learn In-Context? A Case Study of Simple Function
  Classes (2022). \texttt{arXiv:2208.01066}
  \item AttentionViz: A Global View of Transformer Attention (2023). \texttt{arXiv:2305.03210}
  \item Attention Link: An Efficient Attention-Based Low Resource Machine
  Translation Architecture (2023). \texttt{arXiv:2302.00340}
  \item Transformers as Graph-to-Graph Models (2023). \texttt{arXiv:2310.17936}
  \item Stack Attention: Improving the Ability of Transformers to Model
  Hierarchical Patterns (2023). \texttt{arXiv:2310.01749}
  \item Attention Is Not All You Need Anymore (2023). \texttt{arXiv:2308.07661}
  \item LLaMA: Open and Efficient Foundation Language Models (2023). \texttt{arXiv:2302.13971}
  \item GPT-4 Technical Report (2023). \texttt{arXiv:2303.08774}
  \item A Transformer with Stack Attention (2024). \texttt{arXiv:2405.04515}
  \item Conv-Basis: A New Paradigm for Efficient Attention Inference and
  Gradient Computation in Transformers (2024). \texttt{arXiv:2405.05219}
  \item Attention as an RNN (2024). \texttt{arXiv:2405.13956}
  \item On the Power of Convolution Augmented Transformer (2024). \texttt{arXiv:2407.05591}
  \item Memorization in Attention-only Transformers (2024). \texttt{arXiv:2411.10115}
  \item Efficient Attention Mechanisms for Large Language Models: A Survey (2025). \texttt{arXiv:2507.19595}
  \item Forgetting Transformer: Softmax Attention with a Forget Gate (2025). \texttt{arXiv:2503.02130}
  \item Is Random Attention Sufficient for Sequence Modeling? Disentangling Trainable Components in the Transformer (2025). \texttt{arXiv:2506.01115}
  \item Don't Pay Attention (2025). \texttt{arXiv:2506.11305}
  \item Deconstructing Attention: Investigating Design Principles for Effective Language Modeling (2025). \texttt{arXiv:2510.11602}
  \item Understanding Transformers and Attention Mechanisms: An Introduction for Applied Mathematicians (2026). \texttt{arXiv:2604.00965}
  \item 100 Glorious Years of the Ising Model (2026). \texttt{arXiv:2601.11671}
  \item Lectures on singularity categories (2026). \texttt{arXiv:2604.20875}
\end{enumerate}

\end{document}
